% Lösungen Einheit 3 von 5 – Brüche multiplizieren und dividieren (zusammenbau v0.1)
\einheitenkopf{Einheit 3 von 5 · Brüche multiplizieren und dividieren}

\erg{18}{a) $\frac{2}{5} \cdot 35$}
%% TODO Lösung der Erklärzeile 19g) fehlt in der Bank
\erg{19}{a) ja \quad b) $\frac{8}{9}$ \quad c) $\frac{9}{10}$ \quad d) $\frac{10}{11}$ \quad e) $\frac{7}{8}$ \quad f) $\frac{8}{15}$ \quad g) \ldots \quad h) $\frac{6}{35}$ \quad i) $\frac{1 \cdot 1}{2 \cdot 3} = \frac{1}{6}$ \quad j) $\frac{3}{2} \cdot \frac{4}{5} = \frac{12}{10} = \frac{6}{5} = 1\frac{1}{5}$ \quad k) $\frac{7}{15} \cdot \frac{6}{14} \cdot \frac{5}{13} = \frac{210}{2730} = \frac{1}{13}$; $\frac{5}{15} \cdot \frac{4}{14} \cdot \frac{3}{13} = \frac{60}{2730} = \frac{2}{91}$; $210 : 60 = 3{,}5$, also dreieinhalbmal so wahrscheinlich, nicht doppelt. \quad l) $\frac{5}{6} \cdot \frac{1}{5} = \frac{5}{30} = \frac{1}{6}$}
%% TODO Lösung der Erklärzeile 20f) fehlt in der Bank
\erg{20}{a) $\frac{2}{15}$ \quad b) $\frac{5}{21}$ \quad c) $\frac{3}{32}$ \quad d) $\frac{7}{18}$ \quad e) $\frac{4}{15}$ \quad f) \ldots \quad g) $6 \cdot 4 = 24$ \quad h) $\frac{1}{6} \cdot \frac{5}{3} = \frac{5}{18}$ \quad i) $\frac{5}{8} \cdot \frac{16}{15} = \frac{1 \cdot 2}{1 \cdot 3} = \frac{2}{3}$ \quad j) $9 : \frac{3}{4} = 9 \cdot \frac{4}{3} = 12$ Flaschen}
\erg{21}{a) Beim Multiplizieren nimmt man keinen Kehrbruch. Richtig: $\frac{2}{3} \cdot \frac{4}{5} = \frac{8}{15}$. \quad b) 6 Personen sind $\frac{3}{2}$ von 4 Personen: $\frac{3}{2} \cdot \frac{3}{4} = \frac{9}{8} = 1\frac{1}{8}$ Liter \quad c) $\frac{1}{2} \cdot \frac{2}{3} = \frac{2}{6} = \frac{1}{3}$ (4 von 12 Teilen)}

21c)

\bruchrechteck{4}{12}

\erg{22}{a) Er hat den Nenner geteilt, dadurch wird der Bruch größer. Richtig: Nenner mal 4, $\frac{5}{8} : 4 = \frac{5}{32}$. \quad b) In jedes Ganze passen zwei Halbe; in 9 Ganze passen also doppelt so viele Halbe, wie es Ganze sind.}
